\section*{2. 塩化銅水溶液の電気分解}
\begin{enumerate}
    \item 電気分解を行う前の塩化銅水溶液は何色をしているか。（ \textcolor{red}{\textbf{青}} ）色
    \item 塩化銅水溶液に電流を流した。電極A（＋極、陽極）付近で発生する物質の名称と化学式を答えなさい。\\[0.2cm]
    名称（ \textcolor{red}{\textbf{塩素}} ） \hspace{1cm} 化学式（ \textcolor{red}{$\mathbf{Cl_2}$} ）
    \item 電極B（－極、陰極）に付着する物質の名称と化学式を答えなさい。\\[0.2cm]
    名称（ \textcolor{red}{\textbf{銅}} ） \hspace{1cm} 化学式（ \textcolor{red}{$\mathbf{Cu}$} ）
    \item 電気分解を続けると、水溶液の色はどのように変化するか。また、その理由を答えなさい。\\[0.2cm]
    色の変化（ \textcolor{red}{\textbf{青色が薄くなっていく（無色に近づく）。}} ）\\[0.2cm]
    理由（ \textcolor{red}{\textbf{水溶液中の銅イオンが、銅となって電極に付着し減少するから。}} ）
\end{enumerate}
